\documentclass[runningheads]{llncs}

\emergencystretch 1.5em % tolerance for overfull lines; remove for strict typesetting

\usepackage[T1]{fontenc}
\usepackage{graphicx}
\usepackage{epsfig}
\usepackage{subcaption}
\captionsetup{font=footnotesize}
\usepackage{calc}
\usepackage{amsmath}
\usepackage{amssymb}
\usepackage{amstext}
\usepackage{multicol}
%\usepackage{pslatex} % disabled 2026-09-16: Symbol-font bold + \vec accent warnings; default LLNCS (CM) fonts used instead. Re-enable if Times look is wanted.
\usepackage{algorithm2e}
\usepackage[bottom]{footmisc}
\usepackage{booktabs}
\usepackage{multirow}
\usepackage{svg}
\usepackage{hyperref}

\pagenumbering{gobble}
\begin{document}
%
\title{PAMR: Link Prediction on Knowledge Graphs via Path-Aware Multimodal Reasoning}
%
\titlerunning{Link Prediction on KG via Path-Aware Multimodal Reasoning}

\author{Hai-Dang Huynh-Lam\inst{1,2}\orcidID{0000-0001-9487-7281} \and
  Thao Nguyen-Ngoc\inst{1,2} \and
Thanh Le-Ngoc\inst{1,2}\orcidID{0009-0009-6426-2652}}
%
\authorrunning{Hai-Dang et al.}
\institute{
  Faculty of Information Technology, University of Science, Ho Chi Minh City, Vietnam
  \and
  Vietnam National University, Ho Chi Minh City, Vietnam\\
\email{\{hlhdang,nnthao,lnthanh\}@fit.hcmus.edu.vn}}
%
\maketitle  % typeset the header of the contribution
%
\begin{abstract}
  Knowledge graphs organize facts as interconnected entities and relations, yet
  they are always incomplete. Multimodal knowledge graphs supplement structural
  facts with images and text, but prevailing link prediction models treat visual
  and textual data as auxiliary attributes of an attribute-based MMKG (A-MMKG),
  folding them into a single vector per entity. The folding hides the
  connectivity that carries meaning across modalities. We ask what happens when
  the attachments become first-class nodes of a node-based MMKG (N-MMKG). Under
  this modality-as-entity setting, our framework, Path-Aware Multimodal
  Reasoning (PAMR), extracts shortest paths between query entities and scores
  each path with a Transformer over learned knowledge graph embedding
  tokens. On WikiArt-MKG v2,
  over its four evaluation relations, PAMR
  reaches an MRR of 0.379, a 4.4\% gain over the strongest embedding
  baseline, and surpasses it on all four ranking metrics, with Hits@5 at 0.782. Every prediction is accompanied by its extracted multimodal path,
  making the answer traceable to artworks, artists, and shared
  attributes in this experiment.

  \keywords{Knowledge Graph \and Multimodal Learning \and Path-based Reasoning \and Link Prediction \and Transformers.}
\end{abstract}
%
%
%

\section{Introduction}
\label{sec:introduction}

Knowledge graphs (KGs) represent facts as entities wired together by typed relations, and link prediction asks for the connections that are missing. Most methods learn continuous representations from structure alone. Multimodal knowledge graphs (MMKGs) add images and text, and multimodal embedding models do lift link prediction accuracy. Yet prevailing architectures follow the attribute-based MMKG formulation (A-MMKG)~\cite{zhu2024multimodal}, in which visual and textual signals fold into one vector per entity. The folding hides connectivity patterns that cross modality boundaries and leaves reasoning hard to inspect.

Path-aware reasoning goes the opposite way and reads explicit structural sequences between two entities rather than a single proximity score. Yet the models built for this are designed around concept-only graphs, and they tend to tie path modelling to embedding training, a coupling that becomes expensive once multimodal representations are involved. Neither limitation is forced by the data, since the graph can instead follow the node-based MMKG formulation (N-MMKG)~\cite{zhu2024multimodal,liang2024survey}, where multimodal resources become entity nodes in their own right within the relational topology.

We therefore propose the Path-Aware Multimodal Reasoning framework (PAMR). PAMR solves link prediction on multimodal knowledge graphs by scoring structural paths between query entities with a Transformer encoder whose tokens are frozen knowledge graph embeddings. Under this N-MMKG modality-as-entity setting, multimodal information becomes part of the topology itself, so extracted paths can route through shared artworks, styles, or other attributes. The evidence is explicit and traceable, and the decoder-free design sidesteps the initialization cost of continuous sequence modelling.

The core contributions of this work are summarized as follows.

\begin{itemize}
  \item First, we present an architectural framework for path-aware reasoning for link prediction on MMKGs under an N-MMKG modality-as-entity setting, establishing a robust baseline utilizing the unified embeddings provided by VL-KGE.
  \item Second, the proposed path extraction uses both conceptual entities and multimodal entities, with experimental results demonstrating substantial performance gains derived from integrating these multimodal signals.
  \item Third, we introduce a Transformer based sequential architecture for scoring the extracted paths. This approach surpasses the strongest embedding baseline on all four ranking metrics while complementing each prediction with explicit, traceable multimodal reasoning paths, even though the learned Transformer scorer itself is not inherently transparent.
\end{itemize}

\section{Related Work}
\label{sec:related-work}

The PAMR framework sits at the intersection of graph-based relational reasoning, multimodal representation learning, and sequential modelling. We survey each line and state the gap PAMR fills.

\subsection{Path-Based Reasoning on Knowledge Graphs}
\label{sec:rw-path-reasoning}

Path-based reasoning extracts and evaluates multi-hop relational sequences between entities instead of predicting single edges. The lineage starts with statistical random walks such as the Path Ranking Algorithm~\cite{lao2011random}, moves through differentiable logic (NeuralLP~\cite{yang2017differentiable}, RNNLogic~\cite{qu2021rnnlogic}), and reaches reinforcement learning navigators where DeepPath~\cite{xiong2017deeppath} and MINERVA~\cite{das2018go} cast pathfinding as a Markov decision process. Variational inference~\cite{chen2018variational} and hierarchical multi-agent RL~\cite{zhang2022learning} later answered the sparse-reward problem of long paths.

GNN- and sequence-based encoders followed. RED-GNN and NBFNet parameterize the Bellman-Ford recursion with neural components~\cite{zhang2022knowledge,zhu2021neural}, AdaProp prunes the propagation~\cite{zhang2023adaprop}, and SQUIRE translates queries directly into paths with a Transformer~\cite{bai2022squire}. Across all of these, two habits persist: embeddings stay learned rather than frozen, and scoring is driven by message passing, generative decoding, or RL rather than path-sequence Transformers. Moreover, no method is built for multimodal entities, which keeps the embedding-sequence entanglement costly when richer representations enter.

\subsection{Multimodal Knowledge Graph Embeddings}
\label{sec:rw-mmkge}

Multimodal knowledge graph embedding (MMKGE) targets the structural sparsity of KGs
by injecting images and text into representation learning. IKRL~\cite{xie2017image}
first aligned visual and textual features with structural embeddings, and the MMKG collection~\cite{liu2019mmkg} supplied matching baselines. Later systems manage
modality noise and structural complexity: OTE preserves topology through orthogonal
relation transforms~\cite{tang2020orthogonal}, and MMRNS biases negative sampling with
multimodal relations~\cite{xu2022multimodal}.

Transformer-era MMKGE pushes fusion further. MKGformer fuses visual and textual views at multiple levels~\cite{chen2022mkgformer}, MoSE splits modalities before joining them~\cite{zhao2022mose}, VISTA encodes at the triplet level with decoding Transformers~\cite{lee2023vista}, and VL-KGE turns to vision-language models for modality asymmetry at scale~\cite{efthymiou2026vlkge}. Virtually all of these systems operate within the attribute-based paradigm (A-MMKG)~\cite{zhu2024multimodal}, where
multimodal signals condense into a single entity embedding that serves triple scoring, not path discrimination. Node-based graphs (N-MMKGs) do exist, such as Richpedia~\cite{wang2020richpedia}, but their studies centre on multimodal alignment rather than link prediction, and standard multimodal benchmarks such as DB15K bind features to single entities as degree-one leaves, over which path traversal across multimodal nodes is vacuous. PAMR instead targets mediating N-MMKG topologies, where multimodal nodes connect several conceptual entities and mediate multi-hop reasoning paths.

\subsection{Transformers for Knowledge Graph Reasoning}
\label{sec:rw-transformers}

Transformers and pre-trained language models let self-attention contextualise structural data. KG-BERT reads triples as text for sequence classification~\cite{yao2019kg}, and StAR splits contexts to make cross-encoding every candidate cheaper~\cite{wang2021star}. CoKE approaches the graph natively, tokenizing edges and paths into encoder input~\cite{wang2019coke}, while HittER aggregates neighbourhoods hierarchically through masked entity prediction~\cite{chen2021hitter}.

For complex logical queries, KG-Transformer pairs masked pretraining with a sparse Mixture-of-Experts activation~\cite{liu2022mask}. Generative formulations round out the picture: KGT5 casts link prediction as sequence-to-sequence~\cite{saxena2022kgt5}, GenKGC generates targets from relation-guided prompts~\cite{zhu2022genkgc}, and KGTransformer transfers pretrained symbolic structure across KG tasks through prompt tuning~\cite{zhang2023kgtransformer}. What these systems encode, however, are triples, subgraphs, or text, not interleaved multimodal path sequences. They open with language-model pretraining rather than frozen KGE tokens, and they aim at text-only graphs. The mediating N-MMKG topologies targeted by PAMR fall outside this scope as well.

% \subsection{Evaluation Approaches}
% \label{sec:rw-evaluation}
% [Commented draft of negative-sampling survey removed for the SOICT version;
%  see assets/paper.tex for the original text.]

Taken together, the three lines leave open questions that PAMR takes up: how to run path-aware reasoning over heterogeneous multimodal entities, how to project frozen KGE embeddings into a Transformer's token space without an initialization cost, and how to keep sequence ranking stable with a perfectly balanced, query-level paired sampling scheme.

\section{Method}
\label{sec:method}
\subsection{Overview}
\label{sec:overview}

A knowledge graph is a directed multi-relational graph
$\mathcal{G} = (\mathcal{V}, \mathcal{R}, \mathcal{E})$, where $\mathcal{V}$
is the set of entities, $\mathcal{R}$ the set of relation types, and
$\mathcal{E} \subseteq \mathcal{V} \times \mathcal{R} \times \mathcal{V}$ the
set of labelled edges denoting factual triples $(s, r, o)$. In multimodal
knowledge graphs, entities carry perceptual data such as images and text.
We work directly in the node-based formulation (N-MMKG): a multimodal
mapping $\mathcal{M}: \mathcal{V} \rightarrow \mathbb{M}$ with
$\mathbb{M} = \{\text{concept}, \text{image}, \text{text}, \ldots\}$
assigns every entity its modality, and entities with
$\mathcal{M}(v) = \text{concept}$ are conceptual, while every other
entity is a multimodal node that lives in the graph in its own right.
The attribute-based formulation (A-MMKG) instead keeps multimodal
features $\mathcal{X}$ attached to individual entities.

In $\mathcal{G}$, multimodal entities interface with conceptual ones
only through dedicated mediation relations. Conceptual relations join
conceptual entities directly, so any path between distinct conceptual
domains crosses type boundaries solely through shared multimodal hops.
Moreover, one multimodal entity may serve several conceptual entities
at once, which turns shared multimodal assets into first-class bridges
between concepts that per-entity attachment cannot express.

Path-based reasoning on $\mathcal{G}$ treats the structural paths
between query endpoints as the principal evidence for a relation. Given
a query $(s, r, ?)$ with $s$ and $\hat{o}$ conceptual,
$\mathcal{M}(s) = \mathcal{M}(\hat{o}) = \text{concept}$, and $r$ a
relation joining conceptual entities, the goal is to predict $\hat{o}$ by
approximating the conditional distribution $P(\hat{o} \mid s, r,
\mathcal{G})$, which reduces to $P(p_o \mid s, r)$, the distribution
over the paths $p_o$ that connect $s$ and $\hat{o}$ in $\mathcal{G}$. A
learned scoring function $\hat{f}_\theta(p_o, r)$ assigns a real-valued
score to each candidate path, and the entity behind the highest-scoring
path is selected as the answer.

The proposed PAMR consists of four components as illustrated in
Figure~\ref{fig:overview}. First, a frozen pre-trained KGE model supplies
token embeddings for entities and relations. Second, bidirectional BFS
extracts shortest paths between query pairs over the N-MMKG topology.
Third, a Transformer encoder scores each path by contextualising
interleaved entity and relation tokens. Fourth, a margin-based ranking
loss trains the scorer on paired positive and negative paths. Each
component is detailed in the following subsections.

\begin{figure}[t]
  \centering
  \includesvg[width=\linewidth]{figures/overview.svg}
  \caption{Overview of PAMR. Given a query $(s, r, ?)$, frozen KGE embeddings are retrieved for entities and relations. A bidirectional BFS extracts shortest paths between the subject and each candidate object under the modality-as-entity setting. Each path is encoded as an interleaved sequence of entity and relation tokens, prepended with a learnable classification token and the query relation. The Transformer encoder contextualises the sequence, and a two-layer MLP readout produces a path residual added to the frozen KGE triple score. Paired positive and negative paths are scored jointly and trained with a margin-based ranking loss.}
  \label{fig:overview}
\end{figure}

\subsection{Embedding Module}
\label{sec:embedding-module}

To supply pre-encoded relational structure without incurring additional
training cost, we use a frozen pre-trained MMKGE model. Such a model
provides an entity embedding function $g_V: \mathcal{V} \rightarrow
\mathbb{E}$ and a relation embedding function $g_R: \mathcal{R}
\rightarrow \mathbb{E}$ into an embedding space $\mathbb{E}$, together
with a triple scoring function $\mathcal{S}: \mathcal{V} \times
\mathcal{R} \times \mathcal{V} \rightarrow \mathbb{R}$ that assigns
higher values to true triples than to false ones. Both $g_V$ and $g_R$
are learned unified representations that project multimodal entities
and relations, whatever their modality, onto the universal space
$\mathbb{E}$. The embedding functions and the scorer remain frozen
during PAMR training.

\subsection{Path Extraction}
\label{sec:path-extraction}

To provide the Transformer with explicit, interpretable connectivity evidence rather than
relying solely on the implicit proximity captured by KGE embeddings, we extract shortest
paths between query pairs via bidirectional breadth-first search (BFS), bounded by a
maximum hop count $K$.

Formally, a path $p$ between $s$ and $\hat{o}$ comprises $\ell + 1$
entities connected by $\ell$ edges, where
$\ell$ denotes the hop count bounded by $K$.
Paths with $\ell \leq 0$, i.e. unreachable pairs and trivial self-loops,
carry no structural reasoning signal and are excluded from training and
evaluation.

We define two graph configurations that produce qualitatively different
path spaces. The full graph $\mathcal{G}$ keeps every entity type, so
paths may bridge across different types, exploiting cross-type
relationships for reasoning. The concept-only graph $\mathcal{G}_c$
restricts $\mathcal{G}$ to the entities $v$ with
$\mathcal{M}(v) = \text{concept}$ and the edges between them, confining
paths to a homogeneous subgraph that excludes multimodal information.
The latter is the Path-Aware Reasoning (PAR) configuration, the
topology-only counterpart that isolates the contribution of
modality-mediated connectivity.

Under the modality-as-entity formulation, this pairing is the only
coherent control. Multimodal signals and graph topology are entangled by
construction. Multimodal information exists only as explicit nodes, so
any configuration that removes or masks multimodal token features
necessarily also removes the corresponding nodes and their incident
edges. An ablation that retains the full topology while eliminating
multimodal features is therefore not defined in this setting. The one
alternative that preserves connectivity, tokenizing modality nodes with
a constant placeholder, turns every such node into an uninformative
dummy entity, because a modality node's meaning lies precisely in which
concepts it mediates. Such an ablation would measure a corrupted graph
rather than the absence of multimodal features.

Extraction itself follows the evaluation protocol. Every probe masks
the query edge together with its inverse in both orientations before
the search begins, and each split's path graph is built from that
split's own scope, so validation and test facts never enter their own
extraction. The masking is enforced per probe inside the C++ worker
and audited by completeness checks on every output batch.

\subsection{PAMR Architecture}
\label{sec:pamr-architecture}

\subsubsection{Token Interleaving}
\label{sec:token-interleaving}

Each extracted path (Section~\ref{sec:path-extraction}) is encoded as a
flat sequence in the Transformer hidden dimension $d$: entity tokens
carry the embeddings $g_V$ of Section~\ref{sec:embedding-module},
interleaved with edge tokens carrying $g_R$, a learnable classification
token, and the query relation. Let $\mathbf{v}_j \in \mathbb{R}^{d}$ be
the entity token for entity $j$ in the path, obtained from $g_V$.
A path with $\ell$ hops comprises $\ell+1$ entity tokens and $\ell$ edge tokens.
These are interleaved into a token sequence of length $2\ell + 3$:
\begin{equation}
  \tilde{\mathbf{X}} = \bigl[\mathbf{c};\; \mathbf{r};\;
    \mathbf{v}_{s};\; \mathbf{r}_{1};\; \mathbf{v}_{2};\; \mathbf{r}_{2};\;
  \mathbf{v}_{3};\; \ldots;\; \mathbf{v}_{o}\bigr]
  \in \mathbb{R}^{(2\ell+3) \times d}
  \label{eq:token-seq}
\end{equation}
Here $\mathbf{c} \in \mathbb{R}^{d}$ is a learnable classification token
at position 0.
$\mathbf{r}$ is the query relation representation at position 1, while entity tokens
$\mathbf{v}_{s}, \mathbf{v}_{2}, \ldots, \mathbf{v}_{o}$ occupy positions $2 + 2i$
for $i = 0, \ldots, \ell$, and edge tokens $\mathbf{r}_{1}, \ldots, \mathbf{r}_{\ell}$
occupy positions $3 + 2i$ for $i = 0, \ldots, \ell - 1$.
All positions receive absolute positional embeddings
$\mathbf{e}_i = \operatorname{PE}(i)$ via a learned embedding
$\operatorname{PE}: \{0, \ldots, 2K+2\} \to \mathbb{R}^{d}$,
where $K$ is the maximum hop count.
The final token input to the Transformer is $\mathbf{X}_i = \tilde{\mathbf{X}}_i \;+\; \mathbf{e}_i$.
\subsubsection{Transformer Encoder}

\label{sec:transformer-encoder}

The token sequence is contextualised through a stack of $N$
Transformer encoder layers~\cite{vaswani2017attention} with $h$ attention heads each:
\begin{equation}
  \mathbf{H}^{(0)} = \mathbf{X},\qquad
  \mathbf{H}^{(l)} = \operatorname{EncLayer}^{(l)}\bigl(\mathbf{H}^{(l-1)},\; \mathbf{M}\bigr)
  \quad (l = 1, \ldots, N)
  \label{eq:transformer}
\end{equation}
Here $\mathbf{M} \in \{0, 1\}^{2L(\mathcal{B})+1}$ is a padding mask
whose entries $\mathbf{M}_i = 1$ mark invalid positions for paths
shorter than the batch maximum $L(\mathcal{B}) = \max_{p \in \mathcal{B}}
(\ell_p + 1)$ entity tokens.
Each encoder layer applies:
\begin{align}
  \tilde{\mathbf{H}} &= \operatorname{LayerNorm}\bigl(
  \mathbf{H}^{(l-1)} + \operatorname{MHA}(\mathbf{H}^{(l-1)}; \mathbf{M})\bigr)
  \label{eq:attn}\\
  \mathbf{H}^{(l)} &= \operatorname{LayerNorm}\bigl(
  \tilde{\mathbf{H}} + \operatorname{FFN}(\tilde{\mathbf{H}})\bigr)
  \label{eq:ffn}
\end{align}
with $\operatorname{FFN}(\mathbf{x}) = \mathbf{W}^{(2)} \operatorname{GELU}(\mathbf{W}^{(1)} \mathbf{x} + \mathbf{b}^{(1)}) + \mathbf{b}^{(2)}$.
GELU \cite{hendrycks2016gaussian} activation
and Post-LN layer normalisation are used in every encoder layer.

\subsubsection{Readout and Scoring}

\label{sec:readout}

To obtain a fixed-dimensional path representation, the contextualised
token sequence $\mathbf{H}^{(N)} \in \mathbb{R}^{(2\ell+3) \times d}$ is
aggregated by mean pooling over valid tokens. A two-layer MLP maps the
aggregated representation $\mathbf{z}$ to a path residual $\tilde{y}$
with ReLU activation and dropout $p = 0.1$. The path's score adds this
residual to the frozen triple score $\mathcal{S}(s, r, \hat{o})$
(Section~\ref{sec:embedding-module}):
$\hat{y} = \mathcal{S}(s, r, \hat{o}) + \tilde{y}$.
Probes for which extraction finds no path keep the frozen score plus a
learned per-relation offset.

\subsection{Training}

\label{sec:training}

To train a path scorer with stable supervision, we require each query to
contribute equal gradient signal regardless of its number of negative
samples.
\subsubsection{Paired Sampling}
\label{sec:paired-sampling}

For each query $q$, extraction produces one
\textit{positive} path $p_o$, which connects $s$ to the
true object $o$, and $|S^-(q)|$ \textit{negative} paths
$S^-(q) = \{p_{o^-}\}$, which connect
$s$ to sampled false objects $o^-$. The total number of pairs equals
$\sum_q |S^-(q)|$, ensuring that every negative is visited exactly
once per epoch. Collation interleaves positive and negative paths
into flat sequences
$[p_{o,0}, p_{o^-,0}, \ldots, p_{o,B-1}, p_{o^-,B-1}]$.
Here $B$ is the query batch size. Each training batch contains
exactly $B$ positive and $B$ negative samples, producing a guaranteed
balanced composition regardless of the underlying positive-to-negative
ratio. Variable-length paths are zero-padded to the batch maximum
$L(\mathcal{B})$ with the padding masks $\mathbf{M}$ of
Section~\ref{sec:transformer-encoder}.

\subsubsection{Loss Function}
\label{sec:loss-functions}

PAMR is trained with a margin-based ranking loss that groups samples
by $(s, r, ?)$ and enforces a margin between the positive score
and all negative scores for that query. Let $\mathcal{B} = \{(\hat{y}_1, y_1), \ldots, (\hat{y}_{2B}, y_{2B})\}$
denote a training batch with $y_i \in \{0, 1\}$. The loss is defined as:
\begin{equation}
  \mathcal{L}_{\text{rank}} = \frac{1}{|\mathcal{Q}(\mathcal{B})|} \sum_{q \in \mathcal{Q}(\mathcal{B})}\; \frac{1}{|S^+(q)||S^-(q)|} \sum_{\hat{y}^+ \in S^+(q)} \sum_{\hat{y}^- \in S^-(q)}
  \max\bigl(0,\, m - \hat{y}^+ + \hat{y}^-\bigr)
  \label{eq:loss-rank}
\end{equation}
Here $S^{\pm}(q)$ denote the positive and negative path scores of query
$q$ in the batch, each $\hat{y}$ scored as defined in
Section~\ref{sec:readout}, $\mathcal{Q}(\mathcal{B})$ the set of unique
queries in $\mathcal{B}$, and $m = 1.0$ the margin. The per-query
normalisation makes every query's gradient contribution independent of
its negative count, and the margin keeps positive and negative scores
apart.

\section{Experiment}
\label{sec:experiment}
\subsection{Implementation Details}
\label{sec:implementation}

All experiments run on a Vast.ai instance equipped with an NVIDIA
GeForce RTX 5060 Ti and an AMD EPYC 7763 processor allocated with 8 threads
and 16\,GB of virtual RAM. The software stack comprises Python~3.11,
PyTorch~2.7.1 with CUDA~12.8 and cuDNN~9.

\subsubsection{Path Extraction}
The bidirectional BFS is implemented in C++ and called from Python.
Extraction is multi-threaded, and each split's path table is cached on
disk for reuse.

\subsubsection{Training Data and GPU Memory}
The paired dataset of Section~\ref{sec:training} is rebuilt once per epoch
from the precomputed path tables. On the GPU side, the frozen embedding
table dominates memory.

\subsubsection{Training Configuration}
PAMR is trained with AdamW~\cite{loshchilov2019decoupled} at learning
rate $\eta = 0.01$, weight decay $\lambda = 0.01$, $\beta_1 = 0.9$,
$\beta_2 = 0.999$. The ranking loss (Eq.~\ref{eq:loss-rank}) is
used with margin $m = 1.0$. Batch size is 8 queries, i.e. 16 samples
per batch. Training proceeds for 50 epochs. The Transformer encoder uses
$N = 8$ layers and $h = 8$ attention heads with
$d' = 2048$ and dropout $p = 0.1$. The frozen MMKGE model is
VL-KGE~\cite{efthymiou2026vlkge} with TransE embeddings ($D = 768$), and
the token width is $d = 768$. The VL-KGE base model remains frozen
during the training of PAMR.

\subsection{Datasets}
\label{sec:datasets}

Experiments use WikiArt-MKG~v2~\cite{efthymiou2026vlkge}, a multimodal
knowledge graph of fine art with 224,166 entities (216,564 artworks and
7,602 attributes) and 22 relation types, split chronologically into
7,877,220 training, 208,513 validation, and 208,368 test triples.
Artworks carry image embeddings, while styles, genres, years, and locations
carry textual ones, so entity types differ in what modalities they
expose. Training queries get one negative by rejection sampling.
Validation and test queries get ten.

Preprocessing generates inverse
relations for the unidirectional relation types. The similarity
relations \texttt{isRelatedToArtwork} and \texttt{isRelatedToArtist} are
bidirectional and need none. Evaluation targets the four artist-to-artist
relations (\texttt{isInfluencedBy},
\texttt{isInfluencedOn}, \texttt{isPupilOf}, \texttt{isTeacherOf}),
chosen because they can be evaluated consistently under both the
concept-only and the modality-as-entity settings. The comparison
between PAR and PAMR then demonstrates the benefit of the
modality-as-entity formulation over the classic one. To our knowledge, WikiArt-MKG~v2,
published with VL-KGE, is the only public benchmark built this way. In
attachment-formulated graphs such as WN9-IMG or DB15K, each multimodal
feature binds to a single entity, so lifting it into a node would
produce a degree-1 leaf that mediates no path. Modality-as-entity
reasoning would be vacuous there. Porting PAMR to such graphs is
therefore a reformulation rather than a port, and we leave it to
future work.

Path extraction runs bidirectional BFS with at most $K = 5$ hops, base
fanout $F = 100$, and decay factor $\delta = 0.5$~\cite{icaart26}. Both
the full graph $\mathcal{G}$ and the concept-only subgraph $\mathcal{G}_c$
are evaluated, the latter serving as the topology-only
counterpart described in Section~\ref{sec:path-extraction}.

\subsection{Evaluation Metrics}
\label{sec:evaluation-metrics}

For each query, the model assigns a scalar score to every candidate
object in the evaluated candidate set, which holds the gold answer
together with 10 sampled negative candidates as described in
Section~\ref{sec:datasets}, and ranks them in descending order. When
several extracted paths reach one candidate, its score is the maximum
over their path scores, and candidates without an extracted path are
scored as defined in Section~\ref{sec:readout}. We
report Mean Reciprocal Rank (MRR) and Hit-rate at cutoff $k$
(Hits@$k$), overall and per relation type.

\subsection{Results}
\label{sec:results}

Baselines come from VL-KGE~\cite{efthymiou2026vlkge} with four embedding
backbones: TransE~\cite{bordes2013translating}, RotatE~\cite{sun2019rotate}, DistMult~\cite{yang2015embedding}, and ComplEx~\cite{trouillon2016complex}. LLM-based methods are
deliberately excluded since their pretraining on external corpora makes
any comparison unfair. Adapting message-passing path reasoners such as
NBFNet or RED-GNN to the modality-as-entity graph would require
redesigning their message functions around typed modality nodes. We leave
that adaptation to future work and use PAR, the concept-only path scorer,
as the structural baseline.

\subsubsection{Comparative Study}

Table~\ref{tab:results} reports macro-averaged MRR and Hits@$k$ across
the four evaluation relations between artists. PAMR and PAR operate over
the frozen TransE embeddings provided by VL-KGE ($D = 768$), the
translation-based, real-valued backbone that also produces the strongest
baseline. Building on it asks whether path reasoning extracts more from
the same frozen representations than direct triple scoring does. The four
VL-KGE baseline columns come from
our own independent runs of the published VL-KGE scorers, one per
backbone, executed on the same query sets and under the same sampling protocol as PAMR.
They are not the values printed in the VL-KGE paper, which aggregate over
the full relation set. The VL-KGE baselines score triples directly
without path reasoning. PAR and PAMR share the frozen backbone and the
same residual scorer, and they differ only in whether extracted paths
may leave the conceptual subgraph.

\begin{table*}[t]
  \centering
  \setlength{\tabcolsep}{2pt}
  \renewcommand{\arraystretch}{1.15}
  {\scriptsize
    \begin{tabular}{l|cccc|cc}
      \toprule
      \multirow{2}{*}{Metric} &
      \multicolumn{4}{|c|}{VL-KGE Embedding Baselines} &
      \multicolumn{2}{|c|}{Ours} \\
      \cmidrule(lr){2-5} \cmidrule(lr){6-7}
      & TransE & RotatE & DistMult & ComplEx & PAR & PAMR \\
      \midrule
      MRR  & \underline{0.363$\pm$0.030} & 0.327$\pm$0.029 & 0.311$\pm$0.020 & 0.320$\pm$0.048 & 0.221$\pm$0.027 & \textbf{0.379$\pm$0.042} \\
      Hits@1 & \underline{0.345$\pm$0.027} & 0.182$\pm$0.069 & 0.318$\pm$0.076 & 0.283$\pm$0.061 & 0.058$\pm$0.027 & \textbf{0.387$\pm$0.044} \\
      Hits@3 & \underline{0.576$\pm$0.060} & 0.419$\pm$0.098 & 0.528$\pm$0.049 & 0.494$\pm$0.029 & 0.199$\pm$0.041 & \textbf{0.625$\pm$0.044} \\
      Hits@5 & \underline{0.772$\pm$0.065} & 0.597$\pm$0.047 & 0.705$\pm$0.037 & 0.703$\pm$0.062 & 0.331$\pm$0.036 & \textbf{0.782$\pm$0.075} \\
      \bottomrule
  \end{tabular}}
  \caption{Macro-averaged ranking metrics across four evaluation
    relations between artists. \textbf{Bold} indicates the best per metric.
    \underline{Underlines} indicate the second best. Values are reported
    as mean $\pm$ standard deviation. Baseline values are macro-averaged
    over the same four evaluation relations from our independent runs
    and therefore do not correspond to the full relation set reported in
  the VL-KGE publication.}
  \label{tab:results}
\end{table*}

The ordering is consistent across method families: the four VL-KGE
embedding baselines occupy a moderate performance band, PAR falls
below it, and PAMR sits above. TransE leads the baselines with
0.363~MRR, while RotatE, DistMult, and ComplEx stay within
0.02~MRR of one another. Beyond TransE, the choice among the three
scoring functions has limited impact on this graph.

Adding multimodal tokens carries PAMR past the strongest baseline on
all four ranking metrics: MRR rises from 0.363 to 0.379, a 4.4\% gain, with
Hits@1 improving by 12.2\%, Hits@3 by 8.5\%, and Hits@5 by 1.3\%. The
margin is largest at the strict cutoffs and narrows
toward Hits@5, as expected from the residual design, in which the path
residual reorders candidates whose multimodal evidence disagrees with the
frozen triple score. Computationally, the added reasoning is affordable
(Section~\ref{sec:runtime}). Per relation
(Section~\ref{sec:evaluation-metrics}), PAMR leads most metrics on
\texttt{isInfluencedOn} and \texttt{isPupilOf}, the two methods split
\texttt{isInfluencedBy} and \texttt{isTeacherOf}, and PAR trails on
all four.

\subsection{Runtime and Scalability Analysis}
\label{sec:runtime}

We profile the extraction stage on the validation split, whose 202
queries expand to 2,222 candidate probes, one gold answer plus ten
sampled negatives per query (Section~\ref{sec:evaluation-metrics}),
searched by the C++ worker with two CPU threads and own-edge masking
enforced. The worker parses the split graph once and is deterministic
across runs, so single-run wall-clock times are reported.

Loading the graph fixes a constant 3.5--4.6\,s offline baseline at
roughly 0.20\,GB of peak memory. At the production
operating point ($K = 5$, $F = 100$, $\delta = 0.5$), extracting the
complete validation split takes 8.6\,s of wall time, 4.1\,s of which is
search, i.e. 1.8\,ms per probe, with peak memory at 0.23\,GB. The
training split's 15.7M probes scale linearly to roughly eight
CPU-hours at this two-thread rate, a one-time offline cost cached as
path tables and never paid again during training. The recurring cost is
the training loop itself, 543\,s per epoch for PAMR against TransE's
577.

Figure~\ref{fig:bfs-sweep} shows how the two extraction hyperparameters
shape this cost. The hop bound is a structural floor rather than a
free knob. At $K = 2$ no probe yields a path under any fanout or decay. A
creation edge names exactly one artist, so no artwork node connects
two distinct artists, and any artwork-mediated route between two
artists takes at least three hops. The minimum viable bound
is therefore $K = 3$. Reachability then saturates with fanout. At
$K = 5$ coverage rises from 84.5\% $\pm$ 13.4\% at $F = 10$ to 99.4\% at
$F = 50$ and holds there at $F = 100$, with coverage identical across
decay at $F \geq 50$, so the operating point buys a
safety margin rather than extra coverage. The decay factor splits into
the same two regimes. Where fanout is scarce, it gates evidence,
yielding 1{,}542, 1{,}974, and 2{,}115 paths at $K = 5$, $F = 10$ as
$\delta$ runs from 0.2 to 1.0. Once fanout suffices, it only modulates
work, taking the same 2{,}208 paths at $K = 5$, $F = 100$ for
3.8\,s, 8.6\,s, and 14.8\,s at $\delta = 0.2$, 0.5, and 1.0. Extracted
paths average 4.5 hops at the operating point, so the token sequences
of Section~\ref{sec:token-interleaving} exercise their full length
budget of $2K + 3 = 13$, and scoring one candidate costs $O(\ell^{2}h)$
attention over its token sequence.

\begin{figure*}[t]
  \centering
  \includesvg[width=0.98\linewidth]{figures/bfs_sweep.svg}
  \caption{Extraction-stage sensitivity. (a) Share of candidate probes
    with at least one extracted path vs.\ fanout for each hop bound,
    with error bars over the three decay factors. (b) Wall-clock vs
    peak memory at $K = 5$: each line joins the three fanouts at one
  decay factor, the dash style encoding the decay factor and the
  markers denoting the fanout.}
  \label{fig:bfs-sweep}
\end{figure*}

Per-query inference latency grows with the candidate set size. Together these
measurements locate the practical limits of the design: preprocessing scales
with query volume and graph density, scoring scales with the candidate
pool, and both bottlenecks are revisited in
Section~\ref{sec:conclusion}.

\subsection{Case Study}
\label{sec:case-study}

\begin{figure}[t]
  \centering
  \includesvg[width=\linewidth]{figures/case.svg}
  \caption{Case study on the test query (Pablo Picasso,
    \texttt{isInfluencedOn}, ?): candidates Salvador Dali and
    Lee Krasner, a subgraph of intermediate artist relations, and
    the artworks \textit{Man with a Porron} (by Dali),
    \textit{Child with a Dove} (by Picasso), and \textit{Self-Portrait}
  (by Matisse) as multimodal nodes.}
  \label{fig:case-study}
\end{figure}

Figure~\ref{fig:case-study} shows where multimodal entities matter
most. The query (Pablo Picasso, \texttt{isInfluencedOn}, ?) is taken
verbatim from the test set, with two candidates to compare: Salvador
Dali and Lee Krasner. Restricted to artists, PAR's bidirectional BFS
can only walk through colleague relations: Picasso $\xrightarrow{\text{isRelatedToArtist}}$ Juan Gris $\xrightarrow{\text{isRelatedToArtist}}$ Dali, and Picasso $\xrightarrow{\text{isRelatedToArtist}}$ Georges Braque $\xrightarrow{\text{isRelatedToArtist}}$ Matisse $\xleftarrow{\text{isInfluencedBy}}$ Krasner. Both are plausible social routes, but neither carries anything artistic, so the two candidates are hard to tell apart.

With artworks as nodes, PAMR's BFS routes through the creative works. The Dali path re-enters via Picasso's \textit{Child with a Dove} $\xrightarrow{\text{isRelatedToArtwork}}$ \textit{Man with a Porron} $\xrightarrow{\text{hasArtist}}$ Dali. The Krasner path runs through Picasso $\xleftarrow{\text{hasArtist}}$ \textit{Child with a Dove} $\xrightarrow{\text{isRelatedToArtwork}}$ \textit{Self-Portrait} $\xrightarrow{\text{hasArtist}}$ Matisse $\xleftarrow{\text{isInfluencedBy}}$ Krasner.

Given these token sequences, PAMR ranks Dali above Krasner. When the frozen artwork embeddings pass through the encoder, they carry the visual and semantic kinship of the pieces along with them. The Dali route keeps to Spanish avant-garde territory, with shared structural experimentation and foundational Cubist motifs linking Picasso's \textit{Child with a Dove} and Dali's \textit{Man with a Porron}. Reaching Krasner through Matisse's \textit{Self-Portrait}, by contrast, demands a leap from European Cubism through Fauvism to American Abstract Expressionism, movements separated in style, geography, and time, and the path is scored accordingly.

Although the learned Transformer scorer itself is not inherently transparent, its judgement here rests on paths a reader can audit: which artworks the route visits, and what each visit contributes. Compared with the concept-only graph, the reasoning trajectory is grounded in art-historical facts.

\section{Conclusion}
\label{sec:conclusion}

This paper introduced PAMR, a path-aware reasoner for multimodal knowledge graphs that operates under a modality-as-entity setting, in which multimodal values become graph nodes and structural paths are scored by a Transformer over frozen KGE tokens. The framework led all four ranking metrics on WikiArt-MKG v2's four evaluation relations and, more importantly for multimodal practice, attaches to each prediction a path through shared artworks and attributes that a reader can follow end to end. The learned scorer itself remains non-transparent, but what it scores is auditable.

The study has limits. Inference over long interleaved sequences stays
costly when graphs grow. Future work will target this with KV caching
and hardware-aware attention, alongside contrastive objectives that
sharpen separation among hard negatives. Extending the modality-as-entity
formulation to dynamic or temporal modalities is the natural next
setting, and, as traceable evidence becomes a central requirement of
trusted multimodal systems, we consider path-grounded judgement a
strength worth keeping in future designs.

\begin{credits}
  % \subsubsection{\ackname} This research is funded by the University of Science, VNU-HCM, Vietnam under grant number CNTT 2025-16. The authors hereby disclose that an AI tool, namely Google Gemini, has been used to perform literature reviews as well as language and text improvements, and the authors declare that all of the texts generated by the AI have been carefully reviewed manually. No illustration was made with an involvement of any AI tool.

  \subsubsection{\discintname}
  The authors have no competing interests to declare that are
  relevant to the content of this article.
\end{credits}

\bibliographystyle{splncs04}
\bibliography{refs}

\end{document}
